% ================================================================== 3
\section{Fundamentos y aprendizaje autónomo}
\label{sec:fundamentos}

% ------------------------------------------------------------------ 3.1
\subsection{Conceptos y algoritmos requeridos}

\subsubsection{Aritmética modular}

\paragraph{Congruencias.} Dos enteros $a$ y $b$ son congruentes módulo $n$,
$a \equiv b \pmod n$, si $n \mid (a-b)$. Las operaciones de suma y producto
son compatibles con la congruencia, de modo que se puede reducir módulo $n$
después de cada operación sin cambiar el resultado. Esto permite trabajar con
números de tamaño acotado aunque los exponentes sean enormes.

\paragraph{Máximo común divisor e identidad de Bézout.} El algoritmo de
Euclides usa $\gcd(a,b) = \gcd(b, a \bmod b)$ y termina en
$O(\log \min(a,b))$ iteraciones (el peor caso son dos números de Fibonacci
consecutivos, teorema de Lamé)~\cite{knuth,clrs}. La versión extendida
obtiene además enteros $x, y$ con $ax + by = \gcd(a,b)$.

\paragraph{Inverso modular.} $a$ tiene inverso módulo $m$ si y solo si
$\gcd(a, m) = 1$; en ese caso, de $ax + my = 1$ se sigue
$ax \equiv 1 \pmod m$, es decir, $a^{-1} \equiv x \pmod m$.

\paragraph{Teoremas de Fermat y Euler.} Si $p$ es primo y $p \nmid a$,
entonces $a^{p-1} \equiv 1 \pmod p$ (Fermat). En general, si
$\gcd(a,n) = 1$, entonces $a^{\varphi(n)} \equiv 1 \pmod n$ (Euler), donde
$\varphi(n)$ cuenta los enteros en $[1,n]$ coprimos con $n$. Para dos primos
distintos, $\varphi(pq) = (p-1)(q-1)$. La función de Carmichael
$\lambda(pq) = \mathrm{mcm}(p-1, q-1)$ divide a $\varphi(pq)$ y también cumple
$a^{\lambda(n)} \equiv 1 \pmod n$.

\paragraph{Exponenciación binaria.} Escribiendo $e = \sum_i b_i 2^i$ en
binario, $a^e = \prod_{b_i = 1} a^{2^i}$. Basta elevar al cuadrado
$\lfloor\log_2 e\rfloor$ veces y multiplicar una vez por cada bit en 1: el
costo pasa de $O(e)$ a $O(\log e)$ multiplicaciones modulares~\cite{knuth}. En
promedio se hacen $1{,}5$ multiplicaciones por bit del exponente.

\paragraph{Exponenciación por ventana deslizante.} Se precalculan las
potencias impares $a, a^3, \dots, a^{2^k-1}$ y se recorre el exponente de
izquierda a derecha en bloques de hasta $k$ bits que terminan en 1; cada bloque
cuesta una sola multiplicación. El número de multiplicaciones baja a
$\approx b\,(1 + \frac{1}{k+1}) + 2^{k-1}$ para un exponente de $b$ bits, a
cambio de guardar $2^{k-1}$ potencias~\cite{knuth,hac}.

\subsubsection{Pruebas de primalidad}

\paragraph{División de prueba.} Si $n$ es compuesto, tiene un divisor
$d \le \sqrt n$. Basta probar 2, 3 y los candidatos de la forma $6k \pm 1$
(todo primo mayor que 3 tiene esa forma). Es exacta, pero cuesta
$O(\sqrt n) = O(2^{b/2})$ para $n$ de $b$ bits: crece exponencialmente con el
número de bits.

\paragraph{Prueba de Fermat y números de Carmichael.} El teorema de Fermat
sugiere declarar compuesto a $n$ si $a^{n-1} \not\equiv 1 \pmod n$. La prueba
falla con los \emph{números de Carmichael}, compuestos que cumplen
$a^{n-1} \equiv 1$ para \emph{todo} $a$ coprimo con $n$; el menor es
$561 = 3 \cdot 11 \cdot 17$. Por eso no se usa sola.

\paragraph{Miller--Rabin.} Si $p$ es primo impar, las únicas raíces
cuadradas de 1 módulo $p$ son $\pm 1$. Escribiendo $n-1 = 2^s d$ con $d$
impar, para todo $a$ coprimo con un primo $n$ se cumple $a^d \equiv 1$ o
$a^{2^j d} \equiv -1 \pmod n$ para algún $0 \le j < s$. Una base $a$ que
viola esta condición es un \emph{testigo} de que $n$ es
compuesto~\cite{miller1976}. Rabin demostró que, para $n$ compuesto impar, a
lo sumo $1/4$ de las bases no son testigos~\cite{rabin1980}; con $r$ bases
aleatorias independientes, la probabilidad de declarar primo a un compuesto
es como máximo $4^{-r}$. Cada ronda cuesta una exponenciación modular: el
costo total es $O(r\,b\,M(b))$, polinomial en $b$.

\paragraph{Variante determinista.} Con las 12 primeras bases primas
($2, 3, \dots, 37$) la prueba es exacta para todo
$n < 3{,}18 \times 10^{23}$~\cite{sorenson2017}. Para cualquier $n$, el
algoritmo AKS~\cite{aks2004} decide la primalidad en tiempo polinomial
determinista, aunque en la práctica es mucho más lento que Miller--Rabin.

\paragraph{Criba de primos pequeños.} Al generar un primo aleatorio, la
mayoría de los candidatos tiene algún factor primo pequeño. Dividir antes por
los primos menores que una cota $L$ los descarta con divisiones baratas. Por el
teorema de Mertens, la fracción de impares sin factores primos menores que $L$
es $\prod_{3 \le p < L}(1 - 1/p)$; para $L = 2000$ vale $0{,}147$, de modo que
solo el 15\,\% de los candidatos llega a Miller--Rabin.

\subsubsection{Factorización entera}

\paragraph{División de prueba.} Divide sucesivamente entre 2, 3 y $6k \pm 1$
hasta $\sqrt n$. Su peor caso es un semiprimo balanceado $n = pq$ con
$p \approx q \approx \sqrt n$, que es precisamente el tipo de módulo que usa
RSA.

\paragraph{Pollard Rho.} Se itera $x_{i+1} = f(x_i) = x_i^2 + c \bmod n$. Si
$p$ es el menor factor primo de $n$, la sucesión reducida módulo $p$ toma a
lo sumo $p$ valores y, si se comporta como aleatoria, repite un valor tras
$O(\sqrt p)$ pasos esperados por la \emph{paradoja del cumpleaños}. Cuando
$x_i \equiv x_j \pmod p$, $p$ divide a $\gcd(|x_i - x_j|, n)$. La detección
de ciclos de Floyd compara $x_i$ con $x_{2i}$ («tortuga y liebre») usando
memoria constante~\cite{pollard1975}. Como $p \le \sqrt n$, el costo esperado
es $O(n^{1/4})$ iteraciones. Este análisis es \emph{heurístico}: supone que
$f$ se comporta como una función aleatoria.

\paragraph{Variante de Brent.} Brent~\cite{brent1980} detecta el ciclo
comparando $x_i$ con el valor guardado en la última potencia de 2, lo que
requiere una evaluación de $f$ por paso en lugar de tres. Además acumula el
producto de las diferencias $|x - y|$ módulo $n$ y calcula un solo mcd por lote
de pasos. Tiene el mismo orden $O(n^{1/4})$ con una constante menor.

\subsubsection{El criptosistema RSA}

\paragraph{Esquema~\cite{rsa1978}.} Se eligen dos primos distintos $p, q$; se
calcula $n = pq$; se elige $e$ coprimo con $\varphi(n)$ (o con $\lambda(n)$) y
se obtiene $d = e^{-1}$ módulo ese valor. La clave pública es $(n, e)$ y la
privada es $d$. Se cifra con $c = m^e \bmod n$ y se descifra con
$m = c^d \bmod n$, para $0 \le m < n$.

\paragraph{Corrección.} Como $ed = 1 + t\,\varphi(n)$, por Fermat
$m^{ed} \equiv m \pmod p$ y $m^{ed} \equiv m \pmod q$ (también si $p \mid m$),
y por el Teorema Chino del Resto $m^{ed} \equiv m \pmod n$.

\paragraph{Descifrado por el Teorema Chino del Resto.} Con
$d_p = d \bmod (p-1)$, $d_q = d \bmod (q-1)$ y $q_{inv} = q^{-1} \bmod p$ se
calculan $m_1 = c^{d_p} \bmod p$ y $m_2 = c^{d_q} \bmod q$ y se recombinan con
la fórmula de Garner, $m = m_2 + q\,\big(q_{inv}(m_1 - m_2) \bmod p\big)$.
Son dos exponenciaciones con números de la mitad de bits; con multiplicación
escolar cada una cuesta la octava parte, y el total, la cuarta parte del
descifrado directo~\cite{qc1982,hac}.

\paragraph{Seguridad.} Quien factoriza $n$ obtiene $p$, $q$, $\varphi(n)$ y,
con Euclides extendido, $d$: la seguridad de RSA descansa en la dificultad
práctica de factorizar. NIST considera que una clave de 2048 bits ofrece
alrededor de 112 bits de seguridad, y una de 3072, alrededor de
128~\cite{nist80057}. Además, RSA «de libro de texto» es determinista y
maleable; en la práctica se usa con relleno OAEP~\cite{rfc8017}.

\subsubsection{Estándares de RSA que sigue el proyecto}

RSA como algoritmo es uno solo; lo que tiene versiones son los
\textbf{estándares} que establecen cómo usarlo de forma correcta y segura. El
proyecto sigue dos:

\begin{table}[H]
\centering\small
\caption{Estándares aplicados en cada parte del sistema.}
\label{tab:estandares}
\begin{tabularx}{\textwidth}{@{}R{3.0cm} R{3.4cm} Y@{}}
\toprule
\textbf{Parte} & \textbf{Estándar} & \textbf{Qué se toma}\\
\midrule
Generación de claves aleatorias & \textbf{NIST FIPS 186-5}~\cite{fips186}, apéndice A.1 & $e = 65537$; $p$ y $q$
de $b/2$ bits con los dos bits altos en 1; $\gcd(e,p-1)=\gcd(e,q-1)=1$;
$|p-q| > 2^{b/2-100}$; $d = e^{-1} \bmod \lambda(n)$ con $d > 2^{b/2}$.\\
Cifrado y descifrado & \textbf{PKCS \#1 v2.2} (RFC 8017)~\cite{rfc8017} & Primitivas RSAEP
($c = m^e \bmod n$) y RSADP ($m = c^d \bmod n$), incluido el descifrado con la
representación $(p, q, d_P, d_Q, q_{inv})$ de la clave privada (sección 3.2).\\
Relleno & PKCS \#1 v2.2, RSAES-OAEP & No implementado: se usa RSA sin relleno (primitivas RSAEP/RSADP).\\
\bottomrule
\end{tabularx}
\end{table}

Las versiones anteriores de PKCS~\#1 son v1.5 (relleno hoy considerado
vulnerable), v2.0 (introduce OAEP) y v2.1; la v2.2 es la vigente. Los
estándares definen \emph{qué} se calcula y qué condiciones deben cumplir las
claves, no con qué algoritmo hacer la exponenciación ni cómo buscar los primos:
ahí está el espacio de optimización del proyecto. Las claves manuales
(\texttt{generar\_claves(p, q, e)}) usan $\varphi(n)$ para reproducir el
ejemplo resuelto a mano; las aleatorias usan $\lambda(n)$, como exige
FIPS~186-5. FIPS~186-5 exige además claves de al menos 2048 bits; el proyecto
admite claves menores solo con fines educativos.

% ------------------------------------------------------------------ 3.2
\subsection{Fuentes estudiadas y criterios de selección}

Se priorizaron libros de referencia y artículos originales revisados por
pares para los fundamentos, y documentación oficial (normas y documentación
de bibliotecas) para las decisiones técnicas.

\begin{table}[H]
\centering\small
\caption{Evaluación de las fuentes principales.}
\label{tab:fuentes}
\begin{tabularx}{\textwidth}{@{}>{\raggedright}p{4.1cm} >{\raggedright}p{2.4cm} Y@{}}
\toprule
\textbf{Fuente} & \textbf{Tipo} & \textbf{Uso y valoración}\\
\midrule
Cormen et al.~\cite{clrs}, Knuth~\cite{knuth} & Libros de texto & Base de Euclides, exponenciación,
Miller--Rabin y Pollard Rho. Muy confiables y con notación clara; no recogen resultados
recientes sobre bases deterministas.\\
Menezes et al.~\cite{hac} & Manual técnico & Algoritmos criptográficos con pseudocódigo
(ventana deslizante, TCR). Confiable, aunque de 1996: se contrastó con normas actuales.\\
Rivest et al.~\cite{rsa1978}, Miller~\cite{miller1976}, Rabin~\cite{rabin1980},
Pollard~\cite{pollard1975}, Brent~\cite{brent1980}, Quisquater y Couvreur~\cite{qc1982} &
Artículos originales & Fuentes primarias de los algoritmos implementados; permiten citar los
resultados con exactitud.\\
Sorenson y Webster~\cite{sorenson2017}, Agrawal et al.~\cite{aks2004} & Artículos recientes &
Estado del arte en primalidad determinista; base del modo determinista.\\
NIST SP 800-57~\cite{nist80057}, FIPS 186-5~\cite{fips186}, RFC 8017~\cite{rfc8017} & Normas
oficiales & Tamaños de clave, generación de claves y primitivas. Fuente de mayor autoridad para
decisiones de seguridad.\\
Documentación de Python~\cite{pydocs}, pytest~\cite{pytest}, Hypothesis~\cite{hypothesis},
gmpy2~\cite{gmpy2} & Documentación oficial & Configuración de herramientas; se prefirió a
tutoriales no verificados.\\
\bottomrule
\end{tabularx}
\end{table}

\paragraph{Diferencias entre fuentes.} Al comparar las fuentes se
identificaron dos matices que afectan al diseño:
\begin{enumerate}
  \item \textbf{Cota de error de Miller--Rabin.} Los libros de texto dan la
        cota de peor caso $4^{-r}$; para candidatos elegidos al azar el error
        real es mucho menor~\cite{dlp1993}, por eso las normas permiten menos
        rondas en la generación de primos~\cite{fips186}. El proyecto usa
        $r = 20$ por defecto, priorizando la cota de peor caso, porque la
        biblioteca también recibe entradas adversas elegidas por el usuario.
  \item \textbf{Complejidad de Pollard Rho.} Algunas fuentes presentan
        $O(n^{1/4})$ como resultado demostrado; las fuentes primarias aclaran
        que es una estimación heurística. Por eso los experimentos miden
        también iteraciones, no solo tiempo.
\end{enumerate}

\begin{table}[H]
\centering\small
\caption{Estado del arte y papel de cada alternativa en el proyecto. $b$: bits de $n$;
$L_n[\alpha, c] = \exp\!\big((c+o(1))(\ln n)^\alpha(\ln\ln n)^{1-\alpha}\big)$.}
\label{tab:estado-arte}
\begin{tabularx}{\textwidth}{@{}>{\raggedright}p{3.4cm} >{\raggedright}p{2.2cm} >{\raggedright}p{3.0cm} Y@{}}
\toprule
\textbf{Algoritmo} & \textbf{Tipo} & \textbf{Costo} & \textbf{Papel en el proyecto}\\
\midrule
\multicolumn{4}{@{}l}{\emph{Primalidad}}\\
División de prueba & Exacto & $O(2^{b/2})$ & \textbf{Implementado}: primera alternativa.\\
Fermat & Probabilístico & $O(b\,M(b))$ por base & Descartado: falla con Carmichael.\\
Miller--Rabin & Probabilístico & $O(r\,b\,M(b))$ & \textbf{Implementado}: segunda alternativa.\\
MR con 12 bases fijas & Exacto (acotado) & $O(12\,b\,M(b))$ & \textbf{Implementado} (modo determinista).\\
Baillie--PSW & Probabilístico & $O(b\,M(b))$ & No se implementa~\cite{bw1980}.\\
AKS & Exacto & polinomial, muy lento & Solo como referencia teórica~\cite{aks2004}.\\
\midrule
\multicolumn{4}{@{}l}{\emph{Factorización}}\\
División de prueba & Exacto & $O(\sqrt n)$ & \textbf{Implementado}: primera alternativa.\\
Pollard Rho (Floyd) & Aleatorizado & $O(n^{1/4})$ esperado & \textbf{Implementado}.\\
Pollard Rho (Brent) & Aleatorizado & $O(n^{1/4})$, menor constante & \textbf{Implementado}~\cite{brent1980}.\\
Curvas elípticas (ECM) & Aleatorizado & $L_p[1/2, \sqrt 2]$ & Fuera de alcance~\cite{lenstra1987}.\\
Criba del cuerpo de números & Heurístico & $L_n[1/3, 1{,}923]$ & Fuera de alcance; récords públicos (RSA-240, 795 bits)~\cite{boudot2020}.\\
\bottomrule
\end{tabularx}
\end{table}

% ------------------------------------------------------------------ 3.3
\subsection{Conocimientos nuevos que el grupo tuvo que adquirir}

\begin{table}[H]
\centering\footnotesize
\caption{Conocimientos nuevos, por qué fueron necesarios y dónde se aplicaron.}
\label{tab:conocimientos}
\begin{tabularx}{\textwidth}{@{}R{3.4cm} Y Y l@{}}
\toprule
\textbf{Conocimiento} & \textbf{Por qué fue necesario} & \textbf{Dónde se aplicó} & \textbf{Responsable}\\
\midrule
Exponenciación por ventana y ancho óptimo & Reducir multiplicaciones en el descifrado & \texttt{exp\_modular\_ventana}, \texttt{ventana\_optima} & Francesco\\
Aritmética GMP (\texttt{gmpy2}) & Separar costo algorítmico de costo de implementación & \texttt{prototipo\_gmpy2.py}, \texttt{pipeline --gmpy2} & Francesco\\
Testigos fuertes, bases deterministas & Primalidad exacta y rápida bajo $3{,}18\cdot10^{23}$ & \texttt{miller\_rabin\_determinista} & Orlando\\
Criba y teorema de Mertens & Reducir llamadas a Miller--Rabin al generar primos & \texttt{pasa\_criba}, \texttt{generar\_primo} & Orlando\\
pytest e Hypothesis & Pruebas automatizadas y por propiedades & \texttt{tests/} (397 pruebas) & Orlando\\
Detección de ciclos de Brent y lotes de mcd & Factorizar con menos evaluaciones & \texttt{pollard\_rho\_brent} & Camila\\
Ajuste por mínimos cuadrados en escala logarítmica & Contrastar mediciones con la complejidad teórica & \texttt{ajuste\_complejidad.py} & Camila\\
FIPS 186-5 y PKCS \#1 v2.2 & Generar claves y descifrar según los estándares & \texttt{generar\_claves\_\allowbreak aleatorias}, \texttt{descifrar\_crt} & Emmi\\
GitHub Actions & Ejecutar las pruebas en cada \emph{push} en Linux y Windows & \texttt{.github/\allowbreak workflows/\allowbreak pruebas.yml} & Emmi\\
Medición reproducible en Python & Tiempos fiables y conteos exactos & \texttt{benchmark.py} (\texttt{perf\_\allowbreak counter}, \texttt{trace\allowbreak malloc}, \texttt{gc}) & Emmi\\
\bottomrule
\end{tabularx}
\end{table}

% ------------------------------------------------------------------ 3.4
\subsection{Estrategia y evidencias de aprendizaje autónomo}

\paragraph{Estrategia.} Las necesidades de aprendizaje surgen al formular,
diseñar e implementar. Cada necesidad se asigna al integrante responsable del
módulo donde aparece, sin perjuicio de que todos deban comprender el pipeline
completo. Para cada necesidad se define una actividad independiente, una fuente
principal y una \emph{evidencia verificable} en el repositorio (un archivo, una
prueba o un commit propio). El plan se actualiza en cada entrega: la columna
«Estado» indica si la necesidad ya se cubrió y con qué evidencia.

\newcommand{\planhead}{%
\toprule
\textbf{\#} & \textbf{Necesidad} & \textbf{Actividad independiente} & \textbf{Fuente} & \textbf{Evidencia} & \textbf{Estado}\\
\midrule}
\newcolumntype{N}{>{\raggedright\arraybackslash}p{2.9cm}}
\newcolumntype{F}{>{\raggedright\arraybackslash}p{1.5cm}}
\newcolumntype{V}{>{\raggedright\arraybackslash}p{2.9cm}}
\newcolumntype{S}{>{\raggedright\arraybackslash}p{1.5cm}}

\begin{table}[H]
\centering\footnotesize
\caption{Plan de aprendizaje de F.~Canal (aritmética modular).}
\begin{tabularx}{\textwidth}{@{}c N Y F V S@{}}
\planhead
1 & Congruencias, Fermat y Euler & Demostrar a mano la corrección de RSA & \cite{clrs} & Sec.~\ref{sec:fundamentos}, \ref{sec:formulacion} & Cubierto (E1)\\
2 & Invariante de Euclides extendido & Trazar $(17, 3120)$ y comprobar el invariante & \cite{knuth} & Tabla~\ref{tab:traza-euclides}; \texttt{aritmetica.py} & Cubierto (E1)\\
3 & Exponenciación binaria y su costo & Contar multiplicaciones para $e=17$ y $e=65537$ & \cite{knuth} & Tabla~\ref{tab:traza-exp}; prueba A7 & Cubierto (E1)\\
4 & Enteros grandes con \texttt{gmpy2} & Comparar \texttt{powmod} con \texttt{pow} y la implementación propia & \cite{gmpy2} & \texttt{prototipo\_\allowbreak gmpy2.py} & Cubierto (E2)\\
5 & Ventana deslizante y ancho óptimo & Derivar el costo $2^{k-1} + b/(k+1)$ y verificarlo con el contador & \cite{knuth,hac} & \texttt{exp\_modular\_\allowbreak ventana}; A5, A7 & Cubierto (E2)\\
\bottomrule
\end{tabularx}
\end{table}

\begin{table}[H]
\centering\footnotesize
\caption{Plan de aprendizaje de O.~Clemente (primalidad).}
\begin{tabularx}{\textwidth}{@{}c N Y F V S@{}}
\planhead
1 & Testigos fuertes y cota $1/4$ & Leer los artículos originales y trazar $n = 561$ con $a = 2$ & \cite{miller1976,rabin1980} & Tabla~\ref{tab:traza-mr}; P8 & Cubierto (E1)\\
2 & Carmichael y pseudoprimos fuertes & Reunir ejemplos y verificarlos con \texttt{sympy} & \cite{pomerance1980,sorenson2017} & Pruebas P3, P4 & Cubierto (E1)\\
3 & Error de peor caso frente a caso promedio & Contrastar libros y normas & \cite{dlp1993,fips186} & Sec.~\ref{sec:fundamentos} & Cubierto (E1)\\
4 & Pruebas unitarias y por propiedades & Escribir pruebas con \texttt{pytest} e Hypothesis & \cite{pytest,hypothesis} & \texttt{tests/}; X1--X6 & Cubierto (E2)\\
5 & Bases deterministas y criba & Implementar las 12 bases y la criba; medir el ahorro & \cite{sorenson2017} & P9--P11; experimento E5 & Cubierto (E2)\\
\bottomrule
\end{tabularx}
\end{table}

\begin{table}[H]
\centering\footnotesize
\caption{Plan de aprendizaje de C.~Fernández (factorización).}
\begin{tabularx}{\textwidth}{@{}c N Y F V S@{}}
\planhead
1 & Pollard Rho y paradoja del cumpleaños & Leer el artículo original y trazar $n = 3233$ & \cite{pollard1975,clrs} & Tabla~\ref{tab:traza-rho} & Cubierto (E1)\\
2 & Detección de ciclos: Floyd frente a Brent & Comparar ambos en papel e implementarlos & \cite{brent1980,knuth} & \texttt{pollard\_rho\_\allowbreak brent}; F8 & Cubierto (E2)\\
3 & Estado del arte de la factorización & Revisar ECM, cribas y récords & \cite{lenstra1987,boudot2020} & Tabla~\ref{tab:estado-arte} & Cubierto (E1)\\
4 & Límites en algoritmos aleatorizados & Acotar iteraciones y reintentos de Rho & \cite{pollard1975} & \texttt{max\_iter}; F7 & Cubierto (E2)\\
5 & Experimentación y ajuste de modelos & Aprender \texttt{pandas}, \texttt{matplotlib} y mínimos cuadrados & \cite{pydocs} & \texttt{ajuste\_\allowbreak complejidad.py} & Cubierto (E2)\\
\bottomrule
\end{tabularx}
\end{table}

\begin{table}[H]
\centering\footnotesize
\caption{Plan de aprendizaje de E.~Huaman (RSA e infraestructura).}
\begin{tabularx}{\textwidth}{@{}c N Y F V S@{}}
\planhead
1 & Seguridad de RSA y tamaños de clave & Revisar normas y comparar recomendaciones & \cite{nist80057,rfc8017} & Sec.~\ref{sec:introduccion}, \ref{sec:fundamentos} & Cubierto (E1)\\
2 & Git, GitHub y proyectos Python reproducibles & Configurar repositorio, \texttt{src/} y README & \cite{progit,pypa} & Historial; \texttt{README.md} & Cubierto (E1)\\
3 & TCR y recombinación de Garner & Estudiar el descifrado acelerado y verificarlo a mano & \cite{qc1982,hac} & \texttt{descifrar\_crt}; R1, R6 & Cubierto (E2)\\
4 & Generación de claves según FIPS 186-5 & Implementar las condiciones del apéndice A.1 & \cite{fips186} & \texttt{generar\_claves\_\allowbreak aleatorias}; R4, R5 & Cubierto (E2)\\
5 & Integración continua y medición & GitHub Actions; \texttt{perf\_counter}, \texttt{tracemalloc} & \cite{pydocs} & \texttt{pruebas.yml}; \texttt{benchmark.py} & Cubierto (E2)\\
6 & Relleno OAEP & Estudiar RSAES-OAEP y MGF1 & \cite{rfc8017} & -- & En curso\\
\bottomrule
\end{tabularx}
\end{table}

\paragraph{Evidencias.} Cada necesidad cubierta tiene un archivo o una prueba
en el repositorio y un commit hecho desde la cuenta del integrante
(Anexo~\ref{anx:contribuciones}). Las necesidades comunes del grupo
---comprender el pipeline completo, trabajar con Git desde la propia cuenta y
poder ejecutar y explicar cualquier script--- se evidencian en las
demostraciones de los avances (Anexo~\ref{anx:seguimiento}).

